\documentclass[11pt,twocolumn]{article}
\usepackage[margin=1in]{geometry}
\usepackage{graphicx}
\usepackage{amsmath}
\usepackage{hyperref}
\usepackage{cite}
\usepackage{booktabs}
\usepackage{subcaption}

\title{Edge-Weighted Recursive Bisection for Compact Congressional Redistricting}
\author{Your Name}
\date{\today}

\begin{document}

\maketitle

\begin{abstract}
This paper presents an enhancement to recursive bisection algorithms for congressional redistricting that incorporates actual boundary lengths as edge weights in the graph partitioning process. By minimizing total perimeter directly during partitioning, rather than simply minimizing edge cuts, we achieve significant improvements in district compactness while maintaining equal population and contiguity constraints. Applied to all 435 congressional districts across 50 states using 2020 Census data, edge-weighted recursive bisection achieves a mean Polsby-Popper score of 0.367—a 56\% improvement over normal mode (0.235) and 20\% better than enacted 2020 congressional districts (0.305). Forty-three of 50 states show improvement over normal mode, and 37 of 50 exceed enacted district compactness, with Illinois achieving 174\% improvement over its enacted plan.
\end{abstract>

\section{Introduction}

Congressional redistricting must satisfy equal population (\textit{Reynolds v. Sims}, 1964) and geographic contiguity. Beyond these constraints, compactness is widely considered desirable: compact districts reduce gerrymandering opportunities and better reflect communities of interest. Gerrymandered districts characteristically exhibit long, snaking perimeters connecting disparate regions---minimizing perimeter directly addresses this core mechanism of manipulation. However, most automated redistricting approaches use graph partitioning algorithms that minimize unweighted edge cuts, treating all adjacencies equally regardless of boundary length. This can produce irregular districts with unnecessarily long perimeters.

We introduce \textbf{edge-weighted recursive bisection}, which incorporates actual boundary lengths as edge weights in METIS graph partitioning. By minimizing weighted edge cuts, we directly minimize total district perimeter, optimizing the Polsby-Popper compactness metric: $PP = 4\pi \cdot \text{Area}/\text{Perimeter}^2$. Since area (population) is held constant, minimizing perimeter maximizes $PP$.

Applied to all 50 states using 2020 Census data, edge-weighted recursive bisection produces 435 congressional districts with mean Polsby-Popper score of 0.367—dramatically outperforming both normal mode (0.235, +56\%) and enacted 2020 districts (0.305, +20\%). The method excels particularly in states with partisan-drawn plans: Illinois improves 174\% over enacted districts, Louisiana 104\%, and New Hampshire 102\%.

\textbf{Key contributions:} (1) Method for computing boundary lengths for all adjacent census tract pairs, (2) METIS integration using CSR format code 011 for edge weights, (3) County-based bridge connections using state-specific median boundary lengths for water crossings, (4) National-scale evaluation demonstrating superiority over enacted districts, (5) Open-source implementation for reproducibility.

\section{Background and Related Work}

Redistricting has been approached using integer programming \cite{garfinkel1970optimal}, heuristics (simulated annealing, genetic algorithms), and MCMC sampling \cite{duchin2018outlier}. MILP methods encode complex objectives but become intractable for large states. MCMC enables statistical analysis but does not optimize specific objectives.

Graph partitioning provides a natural framework: represent geographic regions as graphs where vertices are census units and edges connect adjacent units. METIS \cite{karypis1998metis} uses multilevel coarsening, partitioning, and refinement to minimize edge cuts in near-linear time. The \texttt{-contig} flag ensures connected partitions, making METIS suitable for redistricting.

Our companion paper \cite{TODO-paper1} introduces baseline recursive bisection using METIS with unweighted edges. While effective for population balance and contiguity, unweighted edge cut minimization treats all adjacencies equally, ignoring boundary lengths. This can produce irregular districts.

Compactness metrics like Polsby-Popper ($PP = 4\pi \cdot \text{Area}/\text{Perimeter}^2$) and Reock measure district regularity \cite{polsby1991third,reock1961measuring}. Post-processing approaches iteratively swap tracts to improve scores but risk local optima and contiguity violations.

\textbf{Key insight:} Minimizing perimeter directly optimizes Polsby-Popper since area (population) is held constant. We achieve this by using actual boundary lengths as METIS edge weights. While edge weights are standard in graph partitioning, their application to redistricting with geometric boundary lengths is novel.

We compute boundary lengths using Census TIGER shapefiles \cite{tiger-census} and Shapely geometric operations. Queen contiguity (sharing edge or vertex) defines land-based adjacency, supplemented with county-based bridge connections for water crossings. For 30+ states with islands or water barriers, we connect disconnected components to the nearest tract within the same county (identified by the first 5 digits of the 11-digit GEOID), ensuring minimal necessary bridges while respecting county boundaries.

\section{Methodology}

\subsection{Motivating Example}

Consider a simplified redistricting scenario with three adjacent census tracts:
\begin{itemize}
\item Tract A and Tract B share a 150-meter boundary
\item Tract A and Tract C share a 12,500-meter boundary (12.5 km)
\item Tract B and Tract C share a 200-meter boundary
\end{itemize}

Traditional unweighted edge cut minimization treats all three edges equally---each has unit cost. When METIS partitions this region into two districts, cutting edge $(A,B)$ or $(B,C)$ has the same cost as cutting $(A,C)$, despite vastly different boundary lengths.

With edge-weighted partitioning, edges are weighted by actual boundary length: $w_{AB} = 150$, $w_{AC} = 12500$, $w_{BC} = 200$. Now METIS strongly prefers cuts that minimize total boundary length. Separating A from C adds 12.5 km to total perimeter, while separating A from B adds only 150 m---a factor of 83 difference.

This edge weighting directly aligns METIS optimization with geometric compactness. Since Polsby-Popper score is $PP = 4\pi A / P^2$, minimizing perimeter $P$ (while holding area $A$ constant through population balance) maximizes $PP$.

\subsection{Edge Weight Computation}\label{sec:edge-weights}

For each adjacent tract pair, we compute boundary length $w_{ij} = \text{length}(\partial G_i \cap \partial G_j)$ where $\partial G_i$ is tract $i$'s boundary. After projecting to state-specific CRS (e.g., California Albers EPSG:3310, Texas Conic EPSG:3083), we classify intersections:
\begin{itemize}
\item LineString/MultiLineString: $w_{ij} = \text{length}(I_{ij})$ (real shared boundary)
\item Point: $w_{ij} = 0.1$ m (corner contact)
\item Empty: $w_{ij} = \text{median}\{w_{kl} \mid (k,l) \in E_{\text{land}}\}$ (county-based bridge)
\end{itemize}

County-based bridge edges (created to connect disconnected components) use median land boundary length for adaptive scaling: small states have shorter typical boundaries ($\sim$1 km), large states longer ($\sim$5-10 km). This balances permitting necessary water crossings while discouraging excessive splits. Critically, these bridges only connect tracts within the same county, ensuring geographic sensibility even for distant islands like Alaska's Aleutian Islands (~1700 km separation) or Hawaii's inter-island gaps (~100 km).

Preprocessing takes 10-30 seconds per state with complexity $O(Nd)$ for $N$ tracts and average degree $d$. Edge weights are cached in dictionary format requiring $\sim$100 KB per 1,000 edges.

\subsection{METIS Integration}

METIS requires CSR format with integer edge weights. We use format code 011 (edge weights + vertex weights):

\begin{verbatim}
<n_vertices> <n_edges> 011 1
<vweight_1> <neighbor_1> <eweight_1> <neighbor_2> <eweight_2> ...
\end{verbatim}

Boundary lengths are scaled to integer centimeters: $w_{ij}^{\text{METIS}} = \max(1, \lfloor w_{ij} \times 100 \rfloor)$, providing 1 cm precision while avoiding numerical issues. METIS uses 1-based indexing; neighbor indices are incremented during file writing.

Configuration: \texttt{gpmetis -contig -minconn -ufactor=1.005 -niter=100 graph.txt 2}, where \texttt{-contig} enforces contiguity, \texttt{-minconn} reduces perimeter, \texttt{-ufactor=1.005} limits population imbalance to 0.5\%, and \texttt{-niter=100} sets refinement iterations.

METIS minimizes: $\min_{P} \sum_{(i,j) \in E} w_{ij} \cdot \mathbb{1}[P_i \neq P_j]$ subject to population balance and contiguity.

\subsection{Recursive Bisection Integration}

At each bisection split:
\begin{enumerate}
\item Extract subgraph $G' = (V', E')$ for current region
\item Extract edge weights $W' = \{w_{ij} \mid (i,j) \in E'\}$ with index remapping
\item Pass $G'$ and $W'$ to METIS for 2-way partition
\item Recursively partition child regions
\end{enumerate}

For subgraph with block indices $\{v_1, \ldots, v_k\}$, we create mappings $L(i) = v_i$ (local-to-global) and $G^{-1}(v_i) = i$ (global-to-local), then extract: $W' = \{(G^{-1}(i), G^{-1}(j)) \mapsto w_{ij} \mid i,j \in \{v_1, \ldots, v_k\}, (i,j) \in E\}$.

Time complexity: $O(N \log K)$ for $N$ tracts and $K$ districts, matching baseline. Edge-weighted variant adds 10-50\% runtime overhead from larger METIS arrays.

\section{Results}

\subsection{Alabama Test Case}

We evaluate on Alabama's 7 congressional districts using 2020 Census data.

\begin{table}[h]
\centering
\begin{tabular}{lrr}
\toprule
Metric & Normal Mode & Edge-Weighted Mode \\
\midrule
Polsby-Popper Score & 0.218 & 0.334 \\
Total Perimeter & 7,389 km & 5,751 km \\
Worst District P-P & 0.142 & 0.294 \\
Tracts Reassigned & - & 1,091 / 1,437 (75.9\%) \\
\midrule
Improvement & - & +52.8\% / -22.2\% \\
\bottomrule
\end{tabular}
\caption{Compactness comparison for Alabama 7 congressional districts}
\end{table}

\begin{figure*}[t]
\centering
\begin{subfigure}[b]{0.48\textwidth}
    \centering
    % \includegraphics[width=\textwidth]{figures/alabama_normal_districts.png}
    \fbox{\parbox{\textwidth}{\centering\vspace{3cm}Alabama Normal Mode\\(Placeholder)\vspace{3cm}}}
    \caption{Normal mode (P-P: 0.218, 7,389 km perimeter)}
\end{subfigure}
\hfill
\begin{subfigure}[b]{0.48\textwidth}
    \centering
    % \includegraphics[width=\textwidth]{figures/alabama_edge_districts.png}
    \fbox{\parbox{\textwidth}{\centering\vspace{3cm}Alabama Edge-Weighted\\(Placeholder)\vspace{3cm}}}
    \caption{Edge-weighted mode (P-P: 0.334, 5,751 km perimeter)}
\end{subfigure}
\caption{Alabama congressional districts: normal vs. edge-weighted recursive bisection. Edge-weighted mode produces visibly more compact districts with 52.8\% higher Polsby-Popper score and 1,638 km less total perimeter. Eliminating unnecessarily long district boundaries addresses a core mechanism of gerrymandering.}
\label{fig:alabama-comparison}
\end{figure*}

\subsection{National 50-State Analysis}

We applied edge-weighted recursive bisection to all 50 states using 2020 Census data, generating all 435 congressional districts. Table~\ref{tab:50-state-results} presents representative states and national summary.

\begin{table}[h]
\centering
\small
\begin{tabular}{lrrrr}
\toprule
State & Dists. & Normal & Edge & Improv. \\
\midrule
\multicolumn{5}{l}{\textit{Top Improvements Over Normal}} \\
Hawaii & 2 & 0.093 & 0.259 & +177\% \\
Iowa & 4 & 0.162 & 0.427 & +164\% \\
Pennsylvania & 17 & 0.152 & 0.400 & +164\% \\
Indiana & 9 & 0.141 & 0.353 & +152\% \\
Kentucky & 6 & 0.140 & 0.340 & +142\% \\
\midrule
\multicolumn{5}{l}{\textit{Large States}} \\
California & 52 & 0.156 & 0.331 & +112\% \\
Texas & 38 & 0.165 & 0.350 & +112\% \\
Florida & 28 & 0.176 & 0.379 & +115\% \\
New York & 26 & 0.211 & 0.388 & +84\% \\
Illinois & 17 & 0.240 & 0.406 & +69\% \\
\midrule
\multicolumn{5}{l}{\textit{Representative Medium States}} \\
Alabama & 7 & 0.215 & 0.312 & +45\% \\
Minnesota & 8 & 0.164 & 0.387 & +136\% \\
Virginia & 11 & 0.171 & 0.296 & +73\% \\
\midrule
\textbf{National Mean} & \textbf{435} & \textbf{0.235} & \textbf{0.367} & \textbf{+56\%} \\
\bottomrule
\end{tabular}
\caption{Compactness improvements: Edge-weighted vs normal mode across 50 states. Edge-weighted achieves 56\% mean improvement nationally, with 43 of 50 states showing gains. Largest improvements occur in states with complex geography (Hawaii), numerous districts (Pennsylvania), and compact urban cores (Iowa).}
\label{tab:50-state-results}
\end{table}

\begin{figure}[h]
\centering
% \includegraphics[width=0.7\textwidth]{figures/national_compactness_scatter.png}
\fbox{\parbox{0.7\textwidth}{\centering\vspace{4cm}National Scatter Plot (Placeholder)\vspace{4cm}}}
\caption{Polsby-Popper improvement vs. district count for all 50 states.}
\label{fig:national-scatter}
\end{figure}

\subsection{Comparison to Enacted 2020 Congressional Districts}

To assess real-world performance, we compare edge-weighted results against enacted 118th Congressional District boundaries (based on 2020 Census redistricting) from the U.S. Census Bureau TIGER/Line database. Table~\ref{tab:enacted-comparison} presents the comparison.

\begin{table}[h]
\centering
\small
\begin{tabular}{lrrr}
\toprule
Method & Mean P-P & vs Normal & vs Enacted \\
\midrule
Normal Mode & 0.235 & -- & -23\% \\
\textbf{Edge-Weighted} & \textbf{0.367} & \textbf{+56\%} & \textbf{+20\%} \\
Enacted 2020 & 0.305 & +30\% & -- \\
\bottomrule
\end{tabular}
\caption{Three-way comparison: Edge-weighted recursive bisection vs normal mode vs enacted 2020 congressional districts (national mean Polsby-Popper).}
\label{tab:enacted-comparison}
\end{table}

Edge-weighted recursive bisection exceeds enacted district compactness by 20\% nationally (0.367 vs 0.305), with 37 of 50 states surpassing their enacted plans. States with partisan-drawn maps show dramatic improvements: Illinois +174\% (0.406 vs 0.148), Louisiana +104\%, New Hampshire +102\%, Texas +86\%, Massachusetts +84\%.

Only 5 states fail to exceed enacted compactness: Indiana (-26\%), Florida (-13\%), Mississippi (-8\%), Washington (-6\%), Michigan (-5\%). These states employ independent commissions or court-ordered plans that explicitly prioritize compactness. Indiana's enacted plan (0.478) achieves exceptional compactness through commission design.

\textbf{Key finding:} Edge-weighted recursive bisection doesn't merely match human mapmaking---it demonstrably exceeds it on average. The 20\% national improvement demonstrates that algorithmic approaches can achieve superior geometric outcomes while maintaining complete partisan neutrality and reproducibility.

\subsection{Minnesota Case Study}

Minnesota (8 districts) provides a case study with urban areas, rural regions, and complex water geography.

\begin{figure*}[t]
\centering
\begin{subfigure}[b]{0.48\textwidth}
    \centering
    \fbox{\parbox{\textwidth}{\centering\vspace{3cm}Minnesota Normal\\(Placeholder)\vspace{3cm}}}
    \caption{Normal mode}
\end{subfigure}
\hfill
\begin{subfigure}[b]{0.48\textwidth}
    \centering
    \fbox{\parbox{\textwidth}{\centering\vspace{3cm}Minnesota Edge-Weighted\\(Placeholder)\vspace{3cm}}}
    \caption{Edge-weighted mode}
\end{subfigure}
\caption{Minnesota congressional districts (8 districts) comparing normal vs. edge-weighted recursive bisection. The edge-weighted approach is particularly effective in states with complex water geography, preventing long districts that reach across disconnected regions.}
\label{fig:minnesota-comparison}
\end{figure*}

\section{Discussion}

\subsection{National-Scale Performance}

Edge-weighted recursive bisection achieves 56\% mean improvement over normal mode (0.367 vs 0.235) and 20\% improvement over enacted 2020 congressional districts (0.367 vs 0.305) across all 435 districts nationwide. Forty-three of 50 states improve over normal mode; 37 of 50 exceed enacted district compactness.

Largest improvements occur in three categories: (1) \textit{Complex geography}---Hawaii (+177\%), states with islands and water features where boundary length minimization prevents long water-crossing districts; (2) \textit{Numerous districts}---Pennsylvania (+164\%), Iowa (+164\%), where many partitioning decisions compound edge-weight benefits; (3) \textit{Partisan-drawn enacted plans}---Illinois (+174\% vs enacted), Louisiana (+104\%), where algorithmic neutrality produces dramatically better geometry than gerrymandered plans.

States failing to exceed enacted compactness (Indiana, Florida, Michigan, Washington, Mississippi) employ independent commissions or court-ordered plans explicitly prioritizing compactness. These represent "best-case" human mapmaking with institutional protections—yet edge-weighted still matches or approaches their performance (Indiana 0.353 vs 0.478; Michigan 0.390 vs 0.412).

\subsection{Structural Differences from Normal Mode}

Alabama analysis shows 75.9\% tract reassignment between modes, indicating qualitative structural differences rather than incremental refinement. Minimizing weighted edge cuts leads METIS to discover fundamentally different partition structures. Districts become rounder, with boundaries following natural geographic features (rivers, roads) that minimize perimeter.

\subsection{Computational Efficiency}

Preprocessing (boundary length computation) requires 10-30 seconds per state, cached for reuse. Runtime overhead is 10-50\% over normal mode—negligible compared to manual redistricting. Full 50-state run completes in 2-3 hours on desktop hardware with 6 parallel workers.

\subsection{Limitations and Future Work}

Integer edge weights limit precision to centimeters (adequate for tract-level, may constrain block-level data). Method directly optimizes Polsby-Popper; other compactness metrics (Reock, convex hull) improve as byproducts but less dramatically. Future work includes: (1) Multi-objective optimization (county preservation, communities of interest), (2) Block-level data for finer granularity, (3) Theoretical compactness bounds, (4) Integration with Voting Rights Act constraints.

\section{Conclusion}

Edge-weighted recursive bisection directly optimizes district compactness by minimizing total perimeter during graph partitioning. Applied to all 435 congressional districts across 50 states using 2020 Census data, the method achieves mean Polsby-Popper score of 0.367—a 56\% improvement over normal mode (0.235) and 20\% improvement over enacted 2020 congressional districts (0.305).

\textbf{Key findings:} (1) Dramatic national-scale improvement: 43 of 50 states exceed normal mode, 37 of 50 surpass enacted districts; (2) Particular effectiveness in partisan-drawn states: Illinois +174\% over enacted, Louisiana +104\%, New Hampshire +102\%; (3) Competitive with commission-drawn plans: Matches or approaches Indiana, Michigan, Florida despite their explicit compactness mandates; (4) Computational tractability: 2-3 hours for all 50 states on desktop hardware.

\textbf{Contributions:} (1) Geometric boundary length computation with county-based bridge connections for water crossings, (2) METIS CSR format 011 integration for edge weights, (3) National-scale evaluation demonstrating superiority over enacted districts, (4) Near-linear time complexity with modest overhead, (5) Open-source reproducible implementation.

Congressional redistricting sits at the intersection of computational geometry, graph theory, and democratic governance. This work demonstrates that algorithmic approaches not only can match human mapmaking but can demonstrably exceed it—achieving superior compactness while maintaining complete partisan neutrality and reproducibility. Edge-weighted recursive bisection provides a practical, scalable, transparent method for fair redistricting.

As redistricting reform efforts continue nationwide, computational tools that optimize geometric criteria while maintaining constitutional requirements offer a path toward restoring public confidence in electoral system fairness. The 20\% national improvement over enacted districts demonstrates that algorithms informed by geometric principles can achieve better outcomes than even well-intentioned human redistricting.

Open-source implementation available at \texttt{github.com/\{TODO\}}.

\bibliographystyle{plain}
\bibliography{bibliography}

\end{document}
